\section{Linear Minimum Distance}

\cite{spielman} gave a proof that any code $\G = \G_1\pi_1\G_2$ where $\G_1,\G2$ are convolutions with constant state size will have sublinear minimum distance. In particular, if ...

 On the otherhand, Benedetto et al. \cite{benedetto2002serial} show that given that $\G_1$ has minmum distance of at least 3, then the minimum distance grows with $n$ as $n^\alpha$ for some $\alpha<1$. We will extend these proof techniques to the base of $\G_1$ having $O(\log n)$ minimum distance and prove that this implies that $\G$ has linear minimum distance. 

 \begin{theorem}
    Let $\G:=\G_1\pi\G_2\in \{0,1\}^{k\times n}$ be a depth-2 serially concatinated code. Let $C_2\in \{0,1\}^{n\times n}$ be a recursive convolution with state size $\sigma_2\leq 1.$ If $\G_1\in\{0,1\}^{k\times n}$ has a state size $\sigma_1=\lceil c\log_2 n\rceil$ for a sufficiently large $c$ and minimum distance $\gamma \sigma_1$ for some constant $\gamma> 0$, then there exists a constant $\delta>0$ such that 
    $$
       \lim_{n\rightarrow\infty} \Pr_{\pi}\left[\frac{\mathsf{minDist}(\G)}{n}\geq \delta\right] = 1
    $$
 \end{theorem}
 \begin{proof}
    For any $v\in \mathsf{im}(\G_1)$, the minimum distance property implies $|v|\geq d_1:=\gamma\sigma_1$.
    We define the expected number of codewords of weight $h$ as $\mathbb{E}[A_h]$. Using the property of the uniform interleaver:
    $$
    \mathbb{E}_{\pi}[A_h] = \sum_{h_1=d_1}^N \frac{A_{h_1}^1 \cdot A_{h_1,h}^2}{\binom{n}{h_1}}
    $$
    where $A^1_h$ is the number of codewords of $G_1$ and with weight $h$, and $A_{h_1,h}^2$ is the number of codewords of $\G_2$ with weight $h$ generated by inputs with weight $h_1$.

    Observe that 
 \end{proof}